\section{1D Solutions of the Schr\"{o}dinger Equation}

We first note that from the time-dependent Schr\"{o}dinger equation, if the potential \(U\pqty{x}\) is finite at some point, then \(\chi\) must be twice-differentiable with a finite second derivative. But if \(U\pqty{x}\) is infinite, then \(\chi\) must be continuous, but its derivative may be discontinuous.

\subsection{Bound States}

\subsubsection{Infinite Potential Well}

Let \(a > 0\). Consider the potential well given by
\begin{equation}
    U\pqty{x} = \begin{cases}
    0, & \qfor \abs{x} \leq a, \\
    \infty, & \quad \text{otherwise}.
    \end{cases}
\end{equation}

Clearly, for \(\abs{x} > a\), we need \(\chi\pqty{x} = 0\), otherwise \(U \chi\) is not finite. We look at solutions to the time-independent Schr\"{o}dinger equation in the region \(\abs{x} \leq a\), which now reads
\begin{equation}
    - \frac{\hbar^2}{2m} \chi'' = E \chi
\end{equation}
and hence
\begin{equation}
    \chi'' + \frac{2 m E}{\hbar^2} \chi = 0.
\end{equation}

Denoting \(k^2 = \frac{2 m E}{\hbar^2}\), we have the general solution
\begin{equation}
    \chi\pqty{x} = A \sin\pqty{k x} + B \cos\pqty{k x}.
\end{equation}

Imposing the boundary conditions of continuity at \(x = \pm a\), we notice that
\begin{equation}
    \begin{split}
        \chi\pqty{a} &= 0 \\
        \chi\pqty{-a} &= 0
    \end{split}
    \quad
    \implies
    \quad
    \begin{split}
        A \sin \pqty{ka} &= 0, \\
        B \cos \pqty{ka} &= 0.
    \end{split}
\end{equation}

For a non-trivial solution where \(\chi \neq 0\), we need one of the following two options to hold.

\begin{enumerate}[label=(\roman*)]
    \item \(A = 0\) and \(\cos \pqty{ka} = 0\), which gives \(k_n = \flatfrac{n \pi}{2 a}\) for odd \(n\) and corresponding eigenfunction \(\chi_n\pqty{X} = B \cos \pqty{k_n x}\).
    \item \(B = 0\) and \(\sin \pqty{ka} = 0\), which gives \(k_n = \flatfrac{n \pi}{2 a}\) for even \(n\) and corresponding eigenfunction \(\chi_n\pqty{X} = A \sin \pqty{k_n x}\).
\end{enumerate}

This allows us to find the allowed energy levels (the eigenvalues of the Hamiltonians) to be
\begin{equation}
    E_n = \frac{\hbar^2 k_n^2}{2 m} = \frac{n^2 \pi^2 \hbar^2}{8 m a^2} \qfor n \in \NN.
\end{equation}

\subsection{Free Particle}

\subsection{Scattering States}